\documentclass[notes=show]{beamer}
%%%%%%%%%%%%%%%%%%%%%%%%%%%%%%%%%%%%%%%%%%%%%%%%%%%%%%%%%%%%%%%%%%%%%%%%%%%%%%%%%%%%%%%%%%%%%%%%%%%%%%%%%%%%%%%%%%%%%%%%%%%%%%%%%%%%%%%%%%%%%%%%%%%%%%%%%%%%%%%%%%%%%%%%%%%%%%%%%%%%%%%%%%%%%%%%%%%%%%%%%%%%%%%%%%%%%%%%%%%%%%%%%%%%%%%%%%%%%%%%%%%%%%%%%%%%
\AtBeginSubsection[
  {\frame<beamer>{\frametitle{Outline}   
    \tableofcontents[currentsection,currentsubsection]}}%
]%
{
  \frame<beamer>{ 
    \frametitle{Outline}   
    \tableofcontents[currentsection,currentsubsection]}
}
\usepackage{ulem}
\usepackage{pdflscape}
\setbeamertemplate{caption}{\insertcaption}
\usepackage{color,xcolor,ucs}
\usepackage{beamerthemesplit}
\usepackage{graphicx}
\usepackage{amsmath}
\usepackage{adjustbox}
\usepackage{pdfpages} 
\definecolor{mintgreen}{rgb}{0.6, 1.0, 0.6}
\makeatletter
\@addtoreset{subfigure}{framenumber}% subfigure counter resets every frame
\makeatother
\newcommand*{\Scale}[2][4]{\scalebox{#1}{$#2$}}%
\newcommand*{\Resize}[2]{\resizebox{#1}{!}{$#2$}}%
\usepackage{amsfonts}
\usepackage{booktabs}
\usepackage{color,soul}
\usepackage{pdfpages}
\setboolean{@twoside}{false}
\usepackage{amsmath}
\usepackage{pdflscape} 
\usepackage{eurosym}
\usepackage{amstext} % for \text
\DeclareRobustCommand{\officialeuro}{%
  \ifmmode\expandafter\text\fi
  {\fontencoding{U}\fontfamily{eurosym}\selectfont e}}
\usepackage[caption = false]{subfig}
\usepackage{appendixnumberbeamer} 
\usepackage{graphicx}
\usepackage{mathpazo}
\usepackage{hyperref}
\usepackage{multimedia}
\usepackage{graphicx}
\usepackage{multirow}
\usepackage{subfig}
\usepackage{verbatim}
\usepackage{booktabs, multicol, multirow}
\setcounter{MaxMatrixCols}{10}
\input{Macros.tex}
\AtBeginSection[]
{
  \begin{frame}[noframenumbering]
    \frametitle{Outline for section \thesection}
    \tableofcontents[currentsection]
  \end{frame}
}


\usepackage{xcolor,soul}
\renewcommand<>{\hl}[1]{\only#2{\beameroriginal{\hl}}{#1}}

% https://tex.stackexchange.com/questions/41683/why-is-it-that-coloring-in-soul-in-beamer-is-not-visible
\makeatletter
\newcommand\SoulColor{%
  \let\set@color\beamerorig@set@color
  \let\reset@color\beamerorig@reset@color}
\makeatother
\SoulColor
\AtBeginSection{\frame{\sectionpage}}

\newsavebox\mybox
\newenvironment{aquote}[1]
  {\savebox\mybox{#1}\begin{quote}\openautoquote\hspace*{-.7ex}}
  {\unskip\closeautoquote\vspace*{1mm}\signed{\usebox\mybox}\end{quote}}

\newtheorem{proposition}[theorem]{Proposition}
\newtheorem{Assumption}[theorem]{Assumption}
\newenvironment{stepenumerate}{\begin{enumerate}[<+->]}{\end{enumerate}}
\newenvironment{stepitemize}{\begin{itemize}[<+->]}{\end{itemize} }
\newenvironment{stepenumeratewithalert}{\begin{enumerate}[<+-| alert@+>]}{\end{enumerate}}
\newenvironment{stepitemizewithalert}{\begin{itemize}[<+-| alert@+>]}{\end{itemize} }
\usetheme{Madrid}
\makeatletter
\setbeamertemplate{footline}
{
  \leavevmode%
  \hbox{%
  \begin{beamercolorbox}[wd=.333333\paperwidth,ht=2.25ex,dp=1ex,center]{author in head/foot}%
    \usebeamerfont{author in head/foot}\insertshortauthor~~\beamer@ifempty{\insertshortinstitute}{}{(\insertshortinstitute)}
  \end{beamercolorbox}%
  \begin{beamercolorbox}[wd=.333333\paperwidth,ht=2.25ex,dp=1ex,center]{title in head/foot}%
    \usebeamerfont{title in head/foot}\insertshorttitle
  \end{beamercolorbox}%
  \begin{beamercolorbox}[wd=.333333\paperwidth,ht=2.25ex,dp=1ex,right]{date in head/foot}%
    \usebeamerfont{date in head/foot}\insertshortdate{}\hspace*{2em}
%    \insertframenumber{} / \inserttotalframenumber\hspace*{2ex} % DELETED
  \end{beamercolorbox}}%
  \vskip0pt%
}
\makeatother
\usepackage{pdfpages}
\begin{document}
	\title[]{\large Labor Demand		\vspace{0.5cm}
	}
	\author[Raffaele Saggio]{Raffaele Saggio}
	\institute[]{UBC}
	\maketitle

%%%%%%%%%%%%%%%%%%%%%%%%%	
%%%%%%%%%%%%%%%%%%%%%%%%%
%%%%%%%%%%%%%%%%%%%%%%%%%
%%%%%%%%%%%%%%%%%%%%%%%%%
\setbeamertemplate{footline}[frame number]{}

%gets rid of bottom navigation symbols
\setbeamertemplate{navigation symbols}{}

%gets rid of footer
%will override 'frame number' instruction above
%comment out to revert to previous/default definitions
\setbeamertemplate{footline}{}
\begin{frame}
\frametitle{Labor Demand}
\begin{itemize}
\item Theory of the firm $\rightarrow$ Theory of labor demand $\rightarrow$ Theory of Inequality
\bigskip
\begin{itemize}
\item Demand for inputs (e.g. labor, capital) from optimization.
\medskip
\item Key question: what explains the elasticity of demand for labor?
\end{itemize}
\bigskip
\item Today: neo-classical labor demand. 
\bigskip 
\item Building block for a variety of ``applied" questions:
\medskip
\begin{itemize}
\medskip
\item Effects of technology
\medskip
\item Effects of immigration / demographics.
\medskip
\item Effects of trade
\medskip
\item ...
\medskip 
\pause
\item  Some of these will populate the ``Inequality" wars chapter. 
\end{itemize}
\end{itemize}
\end{frame}
%%
%%
\begin{frame}
\frametitle{The building block}
\begin{itemize}
\item Firms produce an output combining multiple inputs:
\be
y=f(x_{1},x_{2},\hdots, x_{N})
\ee

\item Typical assumptions: 
\begin{itemize}
\medskip 
\item $f(.)$ exhibits constant return to scale (CRS). 
\medskip
\item work with only two inputs: $y=f(x_{1},x_{2})$.
\end{itemize}

\bigskip
\item Notation time-out
\begin{itemize}
\medskip
\item $f_{i}\equiv$ partial derivative of $f$ with respect to $x_{i}$ with $i\in\{1,2\}$. 
\end{itemize}
\end{itemize}
\end{frame}

\begin{frame}
\frametitle{An Old friend}
\begin{figure}[htb]
\centering
\includegraphics[width=1\textwidth]{aux/isoquanti}
\end{figure}
\end{frame}

\begin{frame}
\frametitle{An Old friend}
\begin{figure}[htb]
\centering
\includegraphics[width=1\textwidth]{aux/isoquanti_aug}
\end{figure}
\end{frame}

\begin{frame}
\frametitle{The Holy Grail: Elasticity of substitution}
\begin{itemize}
\item Under CRS and two inputs: shape of the isoquant is going to be given by
\be
\sigma = - \dfrac{d log(\frac{x_{1}}{x_{2}})}{d log(\frac{f_{1}}{f_{2}})}
\ee
\item Pop-quiz
\pause
\bigskip
\begin{enumerate}
\item Name of the denominator?
\medskip
\item Under perfect competition, what is this denominator?
\end{enumerate}
\bigskip
\pause
\item Name of the game in labor demand (most of the time): estimate $\sigma$ 
\end{itemize}
\end{frame}
%%%%%
\begin{frame}
\frametitle{Special Cases}
\begin{itemize}
\item Cobb-Douglas 
\be
y=f(x_{1},x_{2})= x_{1}^{\alpha}x_{2}^{1-\alpha}
\ee
$\sigma = 1$
\bigskip
\item CES: 
\be
y=f(x_{1},x_{2})=(\alpha x_{1}^{-\rho}+(1-\alpha)x_{2}^{-\rho})^{-\frac{1}{\rho}}
\ee
$\sigma = \dfrac{1}{(1+\rho)}$
\end{itemize}

\end{frame}
%%%%
\begin{frame}
\frametitle{General Case}
\begin{itemize}
\item Want: a general expression for $\sigma$ $\rightarrow$ will be then used to derive ``Marshall (Hicks) Laws"
\pause
\bigskip
\item To make some progress let's consider a cost function $C(w_{1},w_{2},y)$.
\bigskip
\item We can rewrite
\be
C(w_{1},w_{2},y)=\underbrace{\mu(w_{1},w_{2})}_{\text{Unit Cost}}y
\ee
\pause
\item Why?
\pause\
\bigskip
\item Sheppard's Lemma (Envelope Theorem):
\be
x_{1} = C_{1}(w_{1},w_{2},y).
\ee
\be
x_{2} = C_{2}(w_{1},w_{2},y).
\ee
\pause
\item Important: all derivations that are coming up are for a fixed output $y$. 
\end{itemize}
\end{frame}
%%%%
\begin{frame}
\frametitle{General case}
\begin{itemize}
\item We have that $C_{i}$ is homogenous of degree 0 ($HD_{0})$ , then 
\be
x_{1} = C_{1}(\frac{w_{1}}{w_{2}},1,y)\equiv g(\frac{w_{1}}{w_{2}},y)
\ee
\be
x_{2} = C_{2}(\frac{w_{1}}{w_{2}},y)\equiv h(\frac{w_{1}}{w_{2}},y)
\ee
\item Then we have that
\be
C_{11}=\dfrac{\partial x_{1}}{\partial w_{1}} = \dfrac{1}{w_{2}}g_{1}(\frac{w_{1}}{w_{2}},y)
\ee
\be
C_{21}=\dfrac{\partial x_{2}}{\partial w_{1}} = \dfrac{1}{w_{2}}h_{1}(\frac{w_{1}}{w_{2}},y)
\ee
\end{itemize}
\end{frame}
%%%%
\begin{frame}
\frametitle{General Case}
\begin{itemize}
\item Then, for a fixed $y$,
\be
log \frac{x_{1}}{x_{2}} = \log g(\frac{w_{1}}{w_{2}},y) - \log h(\frac{w_{1}}{w_{2}},y) 
\ee
\item Putting the pieces together
\be
\begin{aligned}
\sigma  & = - \dfrac{\left(\frac{w_{1}}{w_{2}}\right) d \log(\frac{x_{1}}{x_{2}})}{d \frac{w_{1}}{w_{2}}} = - \dfrac{w_{1}}{w_{2}}\left(\dfrac{g_{1}(\frac{w_{1}}{w_{2}},y)}{g}-\dfrac{h_{1}(\frac{w_{1}}{w_{2}},y)}{h}\right) \\ 
& = -\dfrac{w_{1}}{w_{2}}\left(\dfrac{w_{2}C_{11}}{g}-\dfrac{C_{21}w_{2}}{h}\right) \\
& \equiv -\dfrac{w_{1}C_{11}}{C_{1}} + \dfrac{w_{1}C_{21}}{C_{2}} 
\end{aligned}
\ee
\end{itemize}
\end{frame}
%%
\begin{frame}
\frametitle{General Case}
\begin{itemize}
\item Now apply Euler's theorem to the function $C_{j}$ ($HD_{0}$)
\pause
\be
C_{11}w_{1} + C_{12}w_{2} = 0 
\ee
\item Based on the above, we arrive at
\be
\sigma =  \dfrac{w_{1}C_{21}}{C_{2}} +  \dfrac{w_{2}C_{12}}{C_{1}} = C_{12}\dfrac{(w_{1}C_{1}+w_{2}C_{2})}{C_{1}C_{2}} 
\ee
\item Applying Euler's theorem to $C(w_{1},w_{2})$ (which recall is a function $HD_{1}$), that is
\be
w_{1}C_{1}+w_{2}C_{2}=C(w_{1},w_{2}) 
\ee
\item We finally obtain
\be
\sigma = \dfrac{C_{12}C}{C_{1}C_{2}}
\ee
\item Extension: partial elasticity of substitution multi-factor production functions
\be
\sigma_{ij} =  \dfrac{C_{ij}C}{C_{i}C_{j}}
\ee
\end{itemize}
\end{frame}
%%
\begin{frame}[shrink=10]
\frametitle{Connecting the dots..}
\begin{itemize}
\item Why these derivations are useful? Because we can start to see the importance of $\sigma$ in driving quantities that we care about...
\be
\begin{aligned}
\dfrac{\partial x_{1}}{\partial w_{2}} & = C_{12} \\
\rightarrow & \epsilon_{12} \equiv \dfrac{w_{2}}{x_{1}}\dfrac{\partial x_{1}}{\partial w_{2}} = \dfrac{w_{2} C_{12}}{x_{1}} \\
& = \dfrac{C_{12}C}{C_{1}C_{2}}\dfrac{w_{2}C_{2}C_{1}}{Cx_{1}} \\ 
& = \sigma \dfrac{w_{2}x_{2}}{C} \\
& = \sigma \underbrace{s_{2}}_{\text{Input 2 cost share}}
\end{aligned}
\ee
\item Also, we have that
\be
\begin{aligned}
& w_{1}\dfrac{\partial x_{1}}{\partial w_{1}} + w_{2}\dfrac{\partial x_{1}}{\partial w_{2}} = 0 \\
& \rightarrow \epsilon_{11} \equiv \dfrac{\partial x_{1}}{\partial w_{1}}\dfrac{w_{1}}{x_{1}} = - \sigma s_{2} = - (1-s_{1})\sigma
\end{aligned}
\ee
\pause
\item Important reminder: we were holding $y$ fixed thus far!
\end{itemize}
\end{frame}
%%
\begin{frame}
\frametitle{Marshall's Laws}
\begin{itemize}
\item We are now ready to derive Marshall's Laws for labor demand.
\bigskip
\item These laws are derived under the following setup:
\begin{itemize}
\bigskip
\item Perfectly competitive \underline{industry} with CRS and cost function $C(w,y)=\mu(w)y$ where $w=(w_{1},\hdots, w_{N})$. 
\bigskip
\item Output is priced at $p=\mu(w)$.
\bigskip
\item Downward sloping demand curve for the industry's output $y=D(p)$.
\end{itemize}
\bigskip 
\item Next, we will show how the elasticity of demand for an input $x_{i}$ relates to three key parameters
\begin{enumerate}
\medskip
\item $\eta\equiv$ elasticity of product demand.
\medskip
\item $\sigma_{ij}\equiv$ partial elasticity of substitution.
\medskip
\item $s_{j}\equiv$ cost shares.
\end{enumerate}
\end{itemize}
\end{frame}
%%
\begin{frame}
\frametitle{Marshall's Laws}
\be
\begin{aligned}
x_{i} & = y\mu_{i}(w) \\
& \rightarrow \log x_{i} = \log y + \log \mu_{i}(w). \\
& \rightarrow \dfrac{\partial \log x_{i}}{\partial w_{j}} = \dfrac{\partial \log y}{\partial \log p}\dfrac{\partial \log p}{\partial w_{j}}+\dfrac{\mu_{ij}(w)}{\mu_{i}(w)} \\
& = -\eta \dfrac{\partial \log \mu(w)}{\partial w_{j}} + \dfrac{\mu_{ij}}{\mu_{i}} \\
& = -\eta \dfrac{\mu_{j}}{\mu} +  \dfrac{\mu_{ij}}{\mu_{i}}\\
& = -\eta \dfrac{\mu_{j}}{\mu} +  \dfrac{\mu_{ij} \mu}{\mu_{i} \mu_{j}} \dfrac{\mu_{j}}{\mu} \\ 
& = -\eta \dfrac{\mu_{j}}{\mu} +  \sigma_{ij}\dfrac{\mu_{j}}{\mu} \\
& = (\sigma_{ij}-\eta)\dfrac{\mu_{j}}{\mu} = (\sigma_{ij}-\eta)\dfrac{x_{j}}{\mu}
\rightarrow w_{j}\dfrac{\partial \log x_{i}}{\partial w_{j}} = (\sigma_{ij}-\eta)\underbrace{\dfrac{w_{j}x_{j}}{\mu}}_{s_{j}} = \epsilon_{ij}-\eta s_{j}
\end{aligned}
\ee
\end{frame}
%%
\begin{frame}
\frametitle{Marshall's Laws}
\begin{itemize}
\item Key relationship:
\be
\dfrac{\partial \log x_{i}}{\partial \log w_{j}}=\underbrace{\epsilon_{ij}}_{\text{Substitution Effect}}-\underbrace{\eta s_{j}}_{\text{Scale Effect}}
\ee
\item \underline{Substitution Effect}: how much does demand for input $x_{i}$ increases for an increase in the price of input $x_{j}$, \textit{holding y fixed}.
\bigskip
\item \underline{Scale Effect}: to understand this, note that since $\mu(w)=p$
\be
\begin{aligned}
\dfrac{\partial \log p}{\partial w_{j}} &  = \dfrac{\mu_{j}(w)}{\mu(w)} \\
& \dfrac{w_{j}\partial \log p}{\partial w_{j}} &  = \dfrac{w_{j}\mu_{j}(w) y}{\mu(w) y} = \dfrac{w_{j}x_{j}}{C} = s_{j}
\end{aligned}
\ee
Thus, when $w_{j}$ increases by 1\%, industry selling price selling rises by $s_{j}$ and this chokes off demand by $\eta s_{j}\%$
\bigskip
\pause
\item Does this bring back any memories?
\end{itemize}
\end{frame}
%%
\begin{frame}
\begin{figure}[htb]
\centering
\includegraphics[width=0.30\textwidth]{aux/Eugen_slutsky_photo}
\end{figure}
\end{frame}
\begin{frame}
\frametitle{Discussion}
\begin{itemize}
\item Neo-classical labor demand framework $\rightarrow$ national or industry wide changes to the wage structure.
\bigskip
\item Firms: all have CRS technology and are price-takers in the input market. 
\bigskip
\item Number and size of the firms is indeterminant.
\bigskip
\pause
\item Very different from other fields!  
\bigskip
\item IO, Trade: lots of firm level heterogeneity in productivity. Does this productivity heterogeneity ``spills" into wages? 
\bigskip
\item Topic of the last two lectures. 
\end{itemize}
\end{frame}
%%
\begin{frame}[shrink=10]
\frametitle{Getting ready for war}
\begin{itemize}
\item Useful benchmark is the Cobb-Douglas CRS benchmark ($\sigma = 1$)
\be
y=f(K,L)=A L^{\alpha} K^{(1-\alpha)}
\ee
\item Note that:
\be
\begin{aligned}
w & = \alpha A L^{\alpha-1}K^{1-\alpha} \\
{wL} & = \alpha A L^{\alpha}K^{1-\alpha} \\
\dfrac{wL}{y} & = \alpha 
\end{aligned}
\ee 
\item $\alpha$ is the labor share (how much of GDP goes to labor?). 
\medskip
\begin{itemize}
\item An entire generation of economists thought that $\alpha=0.66$. 
\item Actually declining over the last years (active area of research).
\end{itemize}
\medskip
\item Starting point to understand productivity trends:
\be
\underbrace{\dfrac{y}{L}}_{\text{Average Product}}=A(\dfrac{K}{L})^{1-\alpha}
\ee
\item Therefore in the C-D marginal product is just average product times a constant
\be
w=\dfrac{\partial y}{\partial L} = \alpha\dfrac{y}{L}
\ee
\end{itemize}
\end{frame}
%%
\begin{frame}
\frametitle{Labor Share is declining in many places...}
\begin{figure}[htb]
\centering
\includegraphics[width=1\textwidth]{aux/labor_share}
\end{figure}
\end{frame}
%%
\begin{frame}
\frametitle{Something deep is happening now....}
\begin{figure}[htb]
\centering
\includegraphics[width=1\textwidth]{aux/productivity_gap}
\end{figure}
\end{frame}
\end{document}

